% !TeX spellcheck = pt_PT

\chapter{Processo de Desenvolvimento e Integração Contínua}
\label{chap:5}

Este capítulo descreve o processo de desenvolvimento adotado no projeto ShopISEL, incluindo a organização dos repositórios, as práticas de controlo de versões, as pipelines de integração e entrega contínua (\ac{CI/CD}) e a estratégia de deployment automatizado.

\section{Organização do Projeto em Repositórios}
\label{sec:5.1}

O ShopISEL adota uma estrutura multi-repositório, em que cada serviço ou componente principal reside num repositório Git independente. Esta abordagem permite que cada equipa ou componente evolua autonomamente, com o seu próprio ciclo de testes, construção e publicação, sem introduzir dependências de build entre serviços distintos.

Os repositórios principais do projeto são:

\begin{itemize}
    \item \texttt{list-service} — serviço de listas de compras;
    \item \texttt{products-service} — serviço de produtos e categorias;
    \item \texttt{prices-service} — serviço de preços;
    \item \texttt{stores-service} — serviço de lojas e georreferenciação;
    \item \texttt{account-service} — serviço de perfil de utilizador;
    \item \texttt{matchqueue-service} — serviço de correspondência e normalização de dados;
    \item \texttt{fcm-notification-worker} — worker de notificações push;
    \item \texttt{scraper-continente} — scraper do website Continente;
    \item \texttt{scraper-pingodoce} — scraper do website Pingo Doce;
    \item \texttt{scraper-orchestrator} — orchestrator dos scrapers via GitHub Actions;
    \item \texttt{frontend} — aplicação web React;
    \item \texttt{frontend-mobile} — aplicação móvel React Native / Expo;
    \item \texttt{deploy} — configuração de infraestrutura e gateway Nginx;
    \item \texttt{keycloak} — configuração do servidor Keycloak;
    \item \texttt{custom-label} — configuração de build e publicação white label da aplicação móvel.
\end{itemize}

\section{Controlo de Versões e Estratégia de Branches}
\label{sec:5.2}

O controlo de versões é assegurado pelo Git, com os repositórios alojados no GitHub. O branch principal de cada repositório é o \texttt{main}, que representa sempre um estado estável e implantável do serviço. O desenvolvimento de novas funcionalidades é realizado em branches separados, que são integrados no \texttt{main} através de Pull Requests após revisão de código.

Esta prática permite que a pipeline de \ac{CI/CD}, configurada no branch \texttt{main}, seja ativada apenas quando o código está pronto para integração, garantindo que apenas código revisto e testado é publicado e implantado.

\section{Pipelines de Integração e Entrega Contínua}
\label{sec:5.3}

Cada repositório de backend inclui uma pipeline de \ac{CI/CD} implementada com GitHub Actions. As pipelines são ativadas automaticamente em cada push para o branch \texttt{main} e seguem, tipicamente, as seguintes etapas:

\begin{enumerate}
    \item \textbf{Build e teste}: o código é compilado e os testes automatizados são executados. Esta etapa garante que o código integrado compila sem erros e que os testes de integração passam antes de avançar para as etapas seguintes.
    \item \textbf{Análise de qualidade}: em alguns serviços, é realizada uma análise estática de código através do SonarCloud, que reporta métricas de qualidade, vulnerabilidades e código duplicado.
    \item \textbf{Construção da imagem Docker}: a imagem Docker do serviço é construída com base no Dockerfile do repositório.
    \item \textbf{Publicação da imagem}: a imagem construída é publicada no GitHub Container Registry (\texttt{ghcr.io/shopisel/<service-name>:latest}), ficando disponível para o processo de deployment.
    \item \textbf{Deploy automático}: as plataformas de hosting (Render.com e Fly.io) são configuradas para detetar novas versões da imagem e realizar o deployment de forma automática.
\end{enumerate}

Para os scrapers, as pipelines incluem uma etapa adicional de execução agendada por cron, que dispara automaticamente o processo de recolha de dados em intervalos regulares, sem necessidade de intervenção manual.

A Figura~\ref{fig:cicd} ilustra o fluxo geral de \ac{CI/CD} para os serviços de backend.

\begin{figure}[H]
\centering
\begin{tikzpicture}[
    node distance=0.5cm and 0.6cm,
    pstep/.style={rectangle, draw, rounded corners=3pt, fill=blue!10,
                 minimum width=2.1cm, minimum height=0.75cm, font=\small, align=center},
    arr/.style={-{Stealth[length=5pt]}, thick}
]
\node[pstep] (push)   {Push para\\main};
\node[pstep, right=0.6cm of push]  (build)  {Build \&\\Testes};
\node[pstep, right=0.6cm of build] (sonar)  {Análise\\SonarCloud};
\node[pstep, right=0.6cm of sonar] (docker) {Build\\Docker};
\node[pstep, right=0.6cm of docker] (reg)   {Publish\\GHCR};
\node[pstep, right=0.6cm of reg]   (deploy) {Deploy\\(Render/Fly)};

\draw[arr] (push.east)   -- (build.west);
\draw[arr] (build.east)  -- (sonar.west);
\draw[arr] (sonar.east)  -- (docker.west);
\draw[arr] (docker.east) -- (reg.west);
\draw[arr] (reg.east)    -- (deploy.west);
\end{tikzpicture}
\caption{Fluxo de \ac{CI/CD} para os serviços de backend}
\label{fig:cicd}
\end{figure}

\section{Reusable Workflows e Orquestração de Scrapers}
\label{sec:5.4}

Um aspeto relevante do processo de desenvolvimento é a utilização de GitHub Actions reusable workflows para a orquestração dos scrapers. Cada repositório de scraper expõe o seu processo de execução como um workflow reutilizável (\texttt{workflow\_call}), que pode ser invocado por repositórios externos.

O repositório \texttt{scraper-orchestrator} contém um workflow que:

\begin{enumerate}
    \item Invoca em paralelo os workflows reutilizáveis de cada scraper;
    \item Aguarda a conclusão de todos os scrapers (\texttt{needs: [...]});
    \item Aciona o MatchQueue Service para processar os dados recolhidos.
\end{enumerate}

Esta arquitetura de orquestração permite adicionar suporte a uma nova insígnia simplesmente criando um novo repositório de scraper com um workflow reutilizável e referenciando-o no orchestrator, sem modificar nenhum dos scrapers existentes.

\section{Ferramentas de Suporte ao Desenvolvimento}
\label{sec:5.5}

O desenvolvimento do ShopISEL recorreu a um conjunto de ferramentas complementares que contribuíram para a qualidade e eficiência do processo:

\begin{itemize}
    \item \textbf{GitHub Copilot}: ferramenta de assistência à escrita de código baseada em modelos de linguagem de grande dimensão (\textit{Large Language Models}), utilizada para sugestões de código, geração de testes e validação de lógica de negócio.
    \item \textbf{SonarCloud}: plataforma de análise estática integrada nas pipelines de \ac{CI/CD}, que reporta métricas de qualidade de código, vulnerabilidades de segurança e cobertura de testes.
    \item \textbf{Docker}: utilizado para containerizar todos os serviços, garantindo reprodutibilidade dos ambientes de desenvolvimento, teste e produção.
    \item \textbf{GitHub Container Registry}: repositório de imagens Docker integrado no GitHub, utilizado para publicar e versionar as imagens de cada serviço.
    \item \textbf{Render.com e Fly.io}: plataformas de cloud hosting utilizadas para o deployment dos serviços, com suporte nativo a imagens Docker e deployment automático por webhook.
\end{itemize}

\section{Gestão de Configuração e Segredos}
\label{sec:5.6}

A configuração de cada serviço é injetada por variáveis de ambiente, seguindo o princípio dos 12-factor apps. Isto permite que a mesma imagem Docker seja executada em diferentes ambientes (desenvolvimento, staging, produção) com configurações distintas, sem necessidade de recompilação.

Os segredos sensíveis (cadeias de ligação a bases de dados, credenciais do Keycloak, credenciais Firebase, tokens de acesso ao MongoDB) são geridos como GitHub Secrets nos repositórios correspondentes, sendo injetados nas pipelines de \ac{CI/CD} e nos ambientes de produção sem exposição no código fonte.

Esta prática é particularmente relevante para os scrapers, cujos workflows de GitHub Actions necessitam de acesso às credenciais do MongoDB e às configurações do MatchQueue Service para funcionar corretamente.

\section{Desenvolvimento White Label}
\label{sec:5.7}

O conceito de White Label aplicado ao ShopISEL refere-se à capacidade de disponibilizar a plataforma sob diferentes identidades visuais e configurações, adaptadas a diferentes contextos de retalho ou parceiros comerciais, sem necessidade de manter bases de código separadas.

A arquitetura multi-repositório e a parametrização das aplicações por variáveis de ambiente criam as condições técnicas para este modelo de distribuição. A aplicação móvel pode ser compilada com diferentes identificadores de aplicação (\textit{bundle ID}), temas visuais, logótipos e apontadores de backend, bastando configurar as variáveis de ambiente e os segredos adequados no repositório de cada variante.

As pipelines de \ac{CI/CD} são o mecanismo central desta abordagem: ao parametrizar o workflow de build com os segredos e configurações específicos de cada parceiro, é possível produzir e publicar uma variante distinta da aplicação sem alterar o código-fonte. A adição de um novo parceiro requer apenas a criação de um conjunto de segredos dedicado e de um workflow que os injete no processo de build existente, reutilizando toda a infraestrutura de testes, construção e deployment já implementada.

O repositório \texttt{custom-label} concretiza esta abordagem, contendo os workflows de GitHub Actions, os ficheiros de configuração Expo e os segredos de build necessários para produzir e publicar uma variante white label da aplicação móvel, com identificador de aplicação, tema visual e apontadores de backend próprios, sem qualquer alteração ao código-fonte partilhado.
